\originalchapter{8.04 作业2：波长与干涉仪}
来源为8.04 Spring 2016 Assignment 2, 2016年2月11日发布, 2月18日17时截止, 共65分. 以下转述全部五题及小问, 解答均为\textbf{AI独立解答}. 题中的加速器预期属于2016年题面背景, 不代表当前建设进度.
\originalsection{第1题：非相对论与相对论波长（20分）}
\begin{exercise}
(a)写电子非相对论波长为$\lambda_{\rm nr}=\delta/\sqrt{E_{\rm kin}}$埃, 能量的数值用eV, 求无量纲$\delta$. (b)对$E=\gamma m_ec^2$, 写$\lambda_{\rm r}=\ell/\sqrt{\gamma^2-1}$, 求$\ell$的fm值和名称. (c)用$\ell,\gamma$重写非相对论式, 比较两波长. (d)计算波长等于电子Compton波长时的能量及速度; 1 TeV电子和7 TeV质子的波长; 非相对论波长相对真实波长误差至多10\%时的最大动能与速度.
\end{exercise}
\begin{solution}
由$p=\sqrt{2m_eE_{\rm kin}}$, $\lambda=h/p$得到$\delta=h/[10^{-10}\sqrt{2m_e(1\,\mathrm{eV})}]=12.264$. 相对论关系$E^2=p^2c^2+m_e^2c^4$给$\ell=h/(m_ec)=2426.31$ fm, 为\textbf{非约化}Compton波长. 这里不是$\hbar/(m_ec)$.

用动能$(\gamma-1)m_ec^2$代入(a), 得$\lambda_{\rm nr}=\ell/\sqrt{2(\gamma-1)}$, 因而$\lambda_{\rm nr}/\lambda_{\rm r}=\sqrt{(\gamma+1)/2}>1$（非零动能）. 若$\lambda=\ell$, 则$p=m_ec$, $\gamma=\sqrt2$, 总能量$722.66$ keV, 动能$211.66$ keV, $\beta=1/\sqrt2=0.7071>0.2$, 显然需要相对论.

高能时$pc\simeq E$, $hc=1.23984\,\mathrm{GeV\,fm}$, 所以1 TeV电子$\lambda\simeq1.23984\times10^{-3}$ fm, 7 TeV质子$\lambda\simeq1.77120\times10^{-4}$ fm, 后者约短7倍. 保留质量项只需用$pc=\sqrt{E^2-m^2c^4}$. 按误差定义$(\lambda_{\rm nr}-\lambda_{\rm r})/\lambda_{\rm r}\leq0.1$, 有$\gamma\leq1.42$, 最大动能$0.42m_ec^2=214.62$ keV, $\beta=\sqrt{1-1/1.42^2}=0.7100$. 若改以近似值作误差分母, 阈值会不同; 本解明确采用真实值为分母.
\end{solution}
\originalsection{第2题：三个电子长度（10分）}
\begin{exercise}
题面采用Gaussian单位$\alpha=e^2/(\hbar c)$、$r_0=e^2/(m_ec^2)$、$\bar\lambda_C=\hbar/(m_ec)$. (a)由$e^2,\hbar,m_e$构造Bohr半径并给fm值. (b)证明三个长度之比为$1:\alpha:\alpha^2$, 并给数值.
\end{exercise}
\begin{solution}
$e^2$的量纲为能量乘长度. 设长度为$\hbar^a m_e^b(e^2)^d$, 比较质量、长度、时间指数得$a=2,b=-1,d=-1$, 故$a_0=\hbar^2/(m_ee^2)=52917.7$ fm. $\bar\lambda_C=\alpha a_0=386.159$ fm, $r_0=\alpha\bar\lambda_C=2.81794$ fm. SI中各$e^2$应替换成$e^2/(4\pi\epsilon_0)$; 这些不同单位制的电荷公式不能混用.
\end{solution}
\originalsection{第3题：概率守恒为何要求幺正（10分）}
\begin{exercise}
双端口归一化入射振幅$u$经过线性矩阵$R$成为$Ru$. 证明对任意入射态保持概率意味着$R^\dagger R=I$.
\end{exercise}
\begin{solution}
条件为$u^\dagger(R^\dagger R-I)u=0$. 令Hermitian矩阵$M=R^\dagger R-I$. 分别取两个单位基矢得$M_{11}=M_{22}=0$; 取$(1,1)/\sqrt2$得$\operatorname{Re}M_{12}=0$, 取$(1,\ii)/\sqrt2$得$\operatorname{Im}M_{12}=0$. 因而$M=0$. 反过来$R^\dagger R=I$直接保证$\norm{Ru}^2=\norm u^2$. 只检验单个输入的概率守恒不能推出幺正.
\end{solution}
\originalsection{第4题：不对称分束器测弹（15分）}
\begin{exercise}
分束器反射、透射概率为$R,T$, $R+T=1$, 振幅$r=\sqrt R,t=\sqrt T$; 矩阵只有右下元素为负. 干涉仪用两个相同分束器, 从上端口输入, 弹在下支路. (a)求幺正矩阵. (b)求坏弹与好弹的两探测概率及爆炸概率. (c)重复测试直到认证或爆炸, 求好弹被无损认证的比例及上界.
\end{exercise}
\begin{solution}
$U=\left(\begin{smallmatrix}r&t\\t&-r\end{smallmatrix}\right)$, $U^\dagger U=U^2=I$. 坏弹不阻断, 输出$U^2(1,0)^{\mathsf T}=(1,0)^{\mathsf T}$, 所以$P_0=1,P_1=0$. 好弹时第一次分束后振幅为$(r,t)$: 下支路概率$T$爆炸, 未爆炸的\textbf{未归一化}振幅为$(r,0)$. 再分束得$(r^2,rt)$, 故$P_{\rm boom}=T$, $P_0=R^2$, $P_1=RT$, 三者和为1. $D_1$认证, $D_0$结果不确定可重试. 几何级数给$f=RT/(1-R^2)=R/(1+R)$（$0<R<1$）, 上确界$1/2$在$R\to1^-$逼近. $R=1$时既不爆炸也不认证, 不能代入说有一半认证成功; 需要的测试次数也随极限增大.
\end{solution}
\originalsection{第5题：平面物质波为何为复指数（10分）}
\begin{exercise}
令$\Psi=\cos\theta+\gamma\sin\theta$, $\theta=kx-\omega t$, 要求任意相移$\epsilon$只乘常数$a(\epsilon)$. 求$\gamma,a$, 指出通常使用的分支.
\end{exercise}
\begin{solution}
逐项比较$\cos\theta,\sin\theta$系数:
$a=\cos\epsilon+\gamma\sin\epsilon$, $a\gamma=-\sin\epsilon+\gamma\cos\epsilon$. 消去$a$得$(1+\gamma^2)\sin\epsilon=0$. 对任意$\epsilon$成立只允许$\gamma=\pm\ii$. 对应$a=\e^{\pm\ii\epsilon}$, $\Psi=\e^{\pm\ii\theta}$. 常规$p=\hbar k,E=\hbar\omega$的正动量正能量平面波选择$+\ii$, 即$\e^{\ii(kx-\omega t)}$; 共轭分支改变导数本征值.
\end{solution}
